% §6.1 Union of Approximate Sets
%
% Cross-reference resolution status:
%
% DEFINED here:
%   \label{sec:union-of-sets}                   -- this section
%   \label{prop:union-fpr-fnr}                  -- the union rate proposition
%
% ALREADY DEFINED (chapter 3):
%   \Cref{ch:boolean-algebra}                   -- chapter 3
%   \Cref{sec:composition-theorem}              -- chapter 3 §3.3
%   \Cref{thm:composition}                      -- chapter 3 §3.3
%   \Cref{sec:or-not}                           -- chapter 3 §3.2
%   \Cref{sec:bridge-to-sets}                   -- chapter 3 §3.5
%
% ALREADY DEFINED (chapter 5):
%   \Cref{ch:approximate-sets}                  -- chapter 5
%   \Cref{sec:from-atom-to-set}                 -- chapter 5 §5.1
%   \Cref{sec:bloom-from-scratch}               -- chapter 5 §5.2
%   \Cref{def:bloom-filter}                     -- chapter 5 §5.2
%   \Cref{sec:one-sided-two-sided}              -- chapter 5 §5.4
%   \Cref{def:one-sided}                        -- chapter 5 §5.4
%
% FORWARD REFS within chapter 6 (resolve in later Plan 6 tasks):
%   \Cref{sec:other-set-ops}                    -- §6.2 (Task 3)
%   \Cref{sec:cancellation}                     -- §6.3 (Task 4)
%   \Cref{sec:perfect-filter-line}              -- §6.4 (Task 5)

\section{Union of Approximate Sets}
\label{sec:union-of-sets}

Chapter~5 framed an approximate set $\ASet{A}$ as a noisy membership
predicate over a universe $\Universe$, with two per-block error rates
$(\fprate, \fnrate)$. The Bloom filter of \Cref{def:bloom-filter} was
the canonical instance: a structurally one-sided approximate set with
$\fnrate = 0$ and a tunable $\fprate$. The present section takes the
first compositional step. Given two approximate sets $\ASet{A}$ and
$\ASet{B}$ over the same universe, we ask what their union is and what
error rates the union inherits. The answer is the composition theorem
of \Cref{thm:composition} applied to the OR gate of
\Cref{sec:or-not}, lifted from the per-bit level to the per-set level.

\subsection*{Definition}

A classical union $A \cup B$ is characterized by its membership
predicate: $\SetIndicator{A \cup B}(x) = \SetIndicator{A}(x) \lor
\SetIndicator{B}(x)$. The set-builder reading is the same:
$x \in A \cup B$ if and only if $x \in A$ or $x \in B$. The membership
test is a Boolean OR of two simpler membership tests. The predicate
framing of \Cref{sec:from-atom-to-set} lifts this characterization
directly to approximate sets.

\begin{definition}[Union of approximate sets]
\label{def:union-of-sets}
Let $\ASet{A}, \ASet{B} : \Universe \to \Bernoulli{\BoolType}$ be two
approximate sets over the same universe $\Universe$. Their
\emph{union} $\ASet{A} \cup \ASet{B}$ is the approximate set whose
membership predicate is the Boolean OR of the two operand predicates:
for every $x \in \Universe$,
\[
  (\ASet{A} \cup \ASet{B})(x) \;=\; \ASet{A}(x) \,\lor\, \ASet{B}(x).
\]
The union approximates the objective set $A \cup B$, where $A$ and $B$
are the underlying objective sets that $\ASet{A}$ and $\ASet{B}$
represent.
\end{definition}

The definition is structural and it makes no commitment to a
particular implementation. To evaluate the union's predicate at $x$,
query $\ASet{A}(x)$ and $\ASet{B}(x)$, then take the OR of the two
Boolean outputs. The union is therefore well-defined as soon as the
two operand predicates are, and its rates will follow from the rates
of the operands together with the OR-gate analysis already done in
\Cref{ch:boolean-algebra}.

\subsection*{Two realizations on Bloom filters: predicate-OR and bit-OR-Bloom}

\Cref{def:union-of-sets} commits to a predicate-level operation:
evaluate $\ASet{A}(x)$, evaluate $\ASet{B}(x)$, take the Boolean OR
of the two outputs. We call this realization \emph{predicate-OR}.
It keeps the two filters as separate objects and asks two queries per
lookup. The proposition derived below (\Cref{prop:union-fpr-fnr}) and
its formula $1 - (1 - \fprate_A)(1 - \fprate_B)$ describe this
construction and only this one.

When both operands are Bloom filters with the same parameters, a
tempting operational shortcut presents itself. Suppose $\ASet{A}$ is
the Bloom filter $(b_1, h_1, \ldots, h_k)$ for a set $A$ and
$\ASet{B}$ is the Bloom filter $(b_2, h_1, \ldots, h_k)$ for a set
$B$, using the same bit-array size $m$ and the same hash family
$h_1, \ldots, h_k$. The bit-wise OR $b_1 \lor b_2$ is itself a
length-$m$ bit array, and one can query the combined bitmap with a
single $k$-position lookup. Call this realization
\emph{bit-OR-Bloom}: one bitmap, one query, one returned verdict.
It is a different operation from predicate-OR, and it gives a
different FPR.

Both realizations preserve FNR-zero on FNR-zero operands. For
bit-OR-Bloom: every element of $A$ has its $k$ bits set in $b_1$,
hence in $b_1 \lor b_2$, and symmetrically for $B$, so every member
of $A \cup B$ is reported $\top$ with certainty. For predicate-OR:
each operand correctly reports $\top$ on its inserted items, and the
OR of any two outputs containing a $\top$ is $\top$.

The two realizations diverge on false positives. Bit-OR-Bloom is
exactly a Bloom filter with parameters $(m, k)$ holding the union of
inserted items, so its FPR is
\begin{equation}
\label{eq:bit-or-bloom-fpr}
  \fprate_{\text{bit-OR}}
  \;\approx\; \bigl(1 - e^{-k(n_A + n_B)/m}\bigr)^k,
\end{equation}
the chapter 5 Bloom-filter formula evaluated at the combined item
count $n_A + n_B$. This is materially different from the predicate-OR
rate $1 - (1 - \fprate_A)(1 - \fprate_B)$. At $m = 5000$, $k = 7$,
$n_A = n_B = 500$, the per-filter rate is $\fprate \approx 0.0082$;
predicate-OR gives $\approx 0.0163$ while
\eqref{eq:bit-or-bloom-fpr} gives $\approx 0.137$, an order of
magnitude higher. The chapter notebook measures both and recovers
each formula within sampling noise.

The gap is structural. Predicate-OR asks
$(\,\forall i: b_1[h_i(x)] = 1\,) \lor (\,\forall i: b_2[h_i(x)] = 1\,)$:
one filter must have \emph{all} $k$ bits set for the query to pass.
Bit-OR-Bloom asks $\forall i: (b_1[h_i(x)] = 1 \lor b_2[h_i(x)] = 1)$:
at every queried position, \emph{either} filter's bit suffices, and
the responsible filter may differ position by position. Queries that
fail predicate-OR can pass bit-OR-Bloom by combining bits from
different sources, so bit-OR-Bloom is the more permissive operation
and carries the larger FPR.

Predicate-OR is the union of approximate sets in the sense of
\Cref{def:union-of-sets}: the lift of $\lor$ at the predicate level,
inheriting its rates from the composition theorem of
\Cref{ch:boolean-algebra}. Bit-OR-Bloom is a different data-structure
choice, the choice of inserting both sets into a single Bloom filter
to save one bitmap and one $k$-position lookup per query, paid for by
the higher FPR of \eqref{eq:bit-or-bloom-fpr}. Bit-OR-Bloom also
requires the two operand filters to share parameters and hash family;
bit-wise OR between filters with different hash families produces a
bit array whose query behavior does not correspond to any clean
set-level operation. Predicate-OR is well-defined regardless.
Exercises in this chapter and the chapter notebook revisit these
distinctions.

\subsection*{False-positive rate}

\begin{proposition}[Union rate formulas]
\label{prop:union-fpr-fnr}
Let $\ASet{A}$ and $\ASet{B}$ be approximate sets of $A, B \subseteq
\Universe$ with per-block rates $(\fprate_A, \fnrate_A)$ and
$(\fprate_B, \fnrate_B)$ respectively. Assume the error events of
$\ASet{A}$ and $\ASet{B}$ at every query point $x \in \Universe$ are
mutually independent. Then the union $\ASet{A} \cup \ASet{B}$ has
false-positive rate
\begin{equation}
\label{eq:union-fpr}
  \fprate_\cup \;=\; 1 - (1 - \fprate_A)(1 - \fprate_B).
\end{equation}
The false-negative rate $\fnrate_\cup$ depends on the prior
distribution over the three regions $A \setminus B$, $B \setminus A$,
$A \cap B$ of the objective union $A \cup B$; for a query distribution
that draws $x$ from each region with respective weights
$w_1, w_2, w_3$ summing to one,
\begin{equation}
\label{eq:union-fnr}
  \fnrate_\cup
  \;=\; w_1 \fnrate_A (1 - \fprate_B)
        \,+\, w_2 \fnrate_B (1 - \fprate_A)
        \,+\, w_3 \fnrate_A \fnrate_B.
\end{equation}
\end{proposition}

We prove the two parts in turn.

\paragraph{FPR derivation.}
Fix a query point $x \notin A \cup B$. Since $x$ is outside both
operand sets, the correct verdict for the union at $x$ is $\bot$, and
the only way for $(\ASet{A} \cup \ASet{B})(x) = \top$ to occur is for
at least one of the two operand predicates to falsely report $\top$
at $x$. The two events $\ASet{A}(x) = \top$ (given $x \notin A$) and
$\ASet{B}(x) = \top$ (given $x \notin B$) are exactly the per-operand
false-positive events at $x$, with conditional probabilities
$\fprate_A$ and $\fprate_B$ respectively. By the cross-set
independence assumption, the two events are independent. We compute
the union false-positive event as the complement of the joint
non-event:
\begin{align*}
  \Prob{(\ASet{A} \cup \ASet{B})(x) = \top \Given x \notin A \cup B}
  &= 1 - \Prob{\ASet{A}(x) = \bot,\, \ASet{B}(x) = \bot
       \Given x \notin A,\, x \notin B} \\
  &= 1 - \Prob{\ASet{A}(x) = \bot \Given x \notin A}\,
         \Prob{\ASet{B}(x) = \bot \Given x \notin B} \\
  &= 1 - (1 - \fprate_A)(1 - \fprate_B).
\end{align*}
The result holds for every $x \notin A \cup B$ with the same
probability, so the union's false-positive rate is
$\fprate_\cup = 1 - (1 - \fprate_A)(1 - \fprate_B)$, as claimed.

\paragraph{FNR derivation.}
A false negative of the union at $x \in A \cup B$ requires both
operand predicates to report $\bot$. Whether each report is itself a
false negative or a true negative depends on which of the three
regions $x$ comes from. Partition the objective union into the
disjoint blocks
\[
  A \cup B \;=\; (A \setminus B) \,\sqcup\, (B \setminus A)
                  \,\sqcup\, (A \cap B).
\]
For $x \in A \setminus B$ we have $x \in A$ and $x \notin B$, so
$\ASet{A}(x) = \bot$ is a false negative of $\ASet{A}$ with
probability $\fnrate_A$, while $\ASet{B}(x) = \bot$ is a true negative
of $\ASet{B}$ with probability $1 - \fprate_B$. By independence the
joint probability that both predicates report $\bot$ on this region is
$\fnrate_A (1 - \fprate_B)$. The region $B \setminus A$ contributes
symmetrically $\fnrate_B (1 - \fprate_A)$. For $x \in A \cap B$ both
operand predicates must miss their respective member, contributing
$\fnrate_A \fnrate_B$ by independence. Weighting the three regions
by the prior $w_1, w_2, w_3$ on $A \cup B$ and applying the law of
total probability yields~\eqref{eq:union-fnr}.

\paragraph{The role of the prior.}
The union false-positive rate $\fprate_\cup$ is a closed-form
polynomial in the two operand FPRs alone; no prior is required because
every negative of $A \cup B$ falls into one and the same region (the
complement of the union). The union false-negative rate, in contrast,
mixes three different per-region rates, and the closed form depends on
how often a query lands in each region. The weights $w_i$ are
determined by the cardinalities $|A \setminus B|, |B \setminus A|,
|A \cap B|$ together with the prior over $A \cup B$. When the operand
sets are disjoint the third region vanishes and the formula
simplifies; when the operands coincide the first two regions vanish
and only the $\fnrate_A \fnrate_B$ term survives. The general
expression is the \emph{marginal} rate over the union; chapter~7's
classification measures will revisit it for confidence-interval
purposes.

\subsection*{Worked example: union of two Bloom filters}

Let $\ASet{A}$ and $\ASet{B}$ be two Bloom filters over the same
universe with the same parameters, with operand rates
\[
  \fprate_A = 0.01, \qquad \fnrate_A = 0,
  \qquad
  \fprate_B = 0.02, \qquad \fnrate_B = 0.
\]
Both operands are FNR-zero, in the sense of \Cref{def:one-sided}: they
sit on the FNR-zero axis of the perfect filter line. The union has
\[
  \fprate_\cup
  \;=\; 1 - (1 - 0.01)(1 - 0.02)
  \;=\; 1 - 0.99 \times 0.98
  \;=\; 1 - 0.9702
  \;=\; 0.0298.
\]
The union FNR is $\fnrate_\cup = 0$, by direct inspection
of~\eqref{eq:union-fnr} with $\fnrate_A = \fnrate_B = 0$: every term
in the weighted sum vanishes. Both realizations of the union
preserve FNR-zero: predicate-OR because the OR of two correct
$\top$'s is $\top$, and bit-OR-Bloom because every member of either
operand set still has all $k$ bits set in the OR'd array. The
arithmetic above, however, is the predicate-OR FPR: it composes the
two operand rates via $1 - (1 - \fprate_A)(1 - \fprate_B)$, the rate
of the Boolean OR of two independent operand verdicts. The bit-OR-Bloom
realization would give a different (larger) number through
\eqref{eq:bit-or-bloom-fpr}, since merging two filters into one
bitmap inflates the load.

The increase from $0.02$ (the worse operand) to $0.0298$ is the price
of the OR composition; the multiplicative second-order term
$\fprate_A \fprate_B = 0.0002$ is small relative to the additive
part, so for small operand rates the union FPR is well-approximated
by the sum $\fprate_A + \fprate_B$.

\subsection*{Connection to the composition theorem}

\eqref{eq:union-fpr} is the composition theorem of
\Cref{thm:composition} applied to the OR gate of \Cref{sec:or-not}
with $k = 2$ stages and $\eta_i = \fprate_i$. Chapter~3 §3.2 already
analyzed the OR-under-noise event at the per-bit level: the false
positive of OR is the failure to read a $\bot$ from at least one
input, and the probability of that joint failure factors into the
product of per-input correctness probabilities $1 - \fprate_i$.
\Cref{sec:bridge-to-sets} previewed the lift to sets; the present
section completes it. The Boolean operation OR is the same at both
levels and the algebra is the same; what changes is the
interpretation of the rates. At the bit level $\fprate_i$ is the
per-bit crossover of an asymmetric channel; at the set level it is the
per-block rate of an approximate set. The composition theorem governs
both.

The connection is the central pedagogical claim of chapters~6 and
later: every set operation has an underlying Boolean gate, every gate
has its noise propagation rule from \Cref{ch:boolean-algebra}, and the
set-level rate algebra is the Boolean noise algebra reread at the
set scale. The remaining sections of this chapter walk through that
correspondence for intersection, complement, and difference, then
identify the regimes in which the composed rates leave the perfect
filter line.

\subsection*{What comes next}

\Cref{sec:other-set-ops} takes up intersection, complement, and
difference. Intersection is the AND gate at the set level, with the
dual structure of the union formulas: multiplicative FPR and an
additive-style FNR. Complement is the NOT gate and swaps the two
rates without further work. Set difference is intersection with the
complemented second operand, and its rates fall out by composing the
two formulas. \Cref{sec:cancellation} then turns to the regimes in
which composition escapes the perfect filter line, when a chain of
operations forces both rates positive even though each operand sat
on an axis, and \Cref{sec:perfect-filter-line} restates the line as
the formal subalgebra it really is.
